\documentclass[11pt,a4paper]{report}
\usepackage{portfolio-informe}
\renewcommand{\contentsname}{Índice}
\renewcommand{\chaptername}{}
\setcounter{secnumdepth}{1}
\setcounter{tocdepth}{1}

\begin{document}

\entregableportada
  {Checklist de resolución}
  {Estado general del portfolio web}
  {Documento de seguimiento independiente\\16 de noviembre de 2026}

\section*{Cómo utilizar este documento}

Cada fila tiene un identificador, prioridad y casilla. Marca la casilla solo después de completar y verificar la tarea.

\begin{itemize}
  \item \textbf{P0:} decisión o bloqueo previo a la entrega/publicación.
  \item \textbf{P1:} funcionamiento correcto, experiencia o riesgo importante.
  \item \textbf{P2:} mejora de calidad/mantenimiento según alcance.
  \item \textbf{DECISIÓN:} necesita acuerdo de la cliente o asesoramiento.
  \item \textbf{VERIFICACIÓN:} no marcar como completada sin evidencia.
\end{itemize}

\tableofcontents
\clearpage

\chapter{Páginas y navegación}

\begin{listaacciones}
  \tarea{P1}{Portada: decidir cuándo aparece la navegación y hacerla accesible también con teclado.}
  \tarea{P1}{Proyectos: corregir el desbordamiento reproducido a 390 px y validar el diseño aprobado.}
  \tarea{P1}{Detalle: decidir y probar el orden de anterior/siguiente y los límites de la lista.}
  \tarea{P1}{Contacto: confirmar una acción funcional para cada instancia de navegación.}
  \tarea{P1}{CV: confirmar si se publicará; si se aprueba, vincular el documento vigente con control accesible.}
  \tarea{P2}{Validar con la cliente nombres, textos, años, imágenes, etiquetas y datos de contacto.}
  \tarea{P2}{Confirmar titularidad, permisos y textos alternativos de las imágenes.}
  \tarea{P1}{Probar apertura/cierre de modales con ratón, teclado y touch.}
\end{listaacciones}
\criteriocierre{Las rutas y acciones acordadas funcionan en las modalidades de entrada seleccionadas y el contenido cuenta con aprobación de la cliente.}

\chapter{Gestión de contenido y CMS}

\begin{listaacciones}
  \tarea{P0}{Acordar quién crea, revisa y publica proyectos y con qué frecuencia.}
  \tarea{P0}{Elegir entre JSON con asistencia técnica, CMS alojado o gestión/API propia; registrar presupuesto y motivos.}
  \tarea{P1}{Aprobar campos, bloques, idiomas, límites y validaciones.}
  \tarea{P1}{Definir almacenamiento de imágenes, formatos, compresión, texto alternativo y derechos.}
  \tarea{P1}{Definir usuarios, roles, autenticación y permisos si existe panel/API.}
  \tarea{P1}{Probar crear, previsualizar, publicar, editar, retirar y recuperar contenido.}
  \tarea{P2}{Decidir el destino de los cinco proyectos actuales y validar su migración, si aplica.}
  \tarea{P2}{Aprobar \archivo{double-image}: implementar su presentación o retirar el tipo no soportado.}
\end{listaacciones}
\criteriocierre{La usuaria completa el flujo editorial aprobado sin modificar código y existe una recuperación probada.}

\chapter{Licencia, autoría y derechos}

\begin{listaacciones}
  \tarea{P0}{Identificar titulares/autorizaciones de código, diseño, textos, imágenes, fuentes, marca y contribuciones.}
  \tarea{P0}{Definir usos permitidos y comerciales, autorización, pago y atribución requerida.}
  \tarea{P0}{Revisar/redactar la licencia con asesoramiento legal; contrastar MIT actual y el aviso de copyright.}
  \tarea{P0}{Separar condiciones de terceros que no puedan relicenciarse junto con el proyecto.}
  \tarea{P1}{Actualizar avisos únicamente tras aprobación de titulares y condiciones.}
\end{listaacciones}
\criteriocierre{La licencia y los avisos corresponden a derechos acreditados; no quedan materiales de terceros asumidos como propiedad de la cliente.}

\chapter{Rutas, errores y estados degradados}

\begin{listaacciones}
  \tarea{P1}{Añadir una respuesta de ruta 404 con navegación de salida.}
  \tarea{P1}{Definir recuperación ante errores de renderizado y carga, con mensajes útiles.}
  \tarea{P1}{Probar ID no numérico, inexistente, primero y último; hacer consistente la navegación.}
  \tarea{P1}{Verificar acceso directo y recarga de rutas en el alojamiento de producción.}
  \tarea{P1}{Probar imágenes, vídeos y fuentes inaccesibles sin bloquear el resto de la página.}
  \tarea{P2}{Actualizar título, favicon e idioma del documento.}
\end{listaacciones}
\criteriocierre{Rutas válidas, desconocidas y fallidas tienen una respuesta deliberada en desarrollo y producción.}

\chapter{Experiencia, accesibilidad y adaptación}

\begin{listaacciones}
  \tarea{P1}{Resolver el overflow de \rutaweb{/projects}: documento medido en 1160 px con viewport de 390 px.}
  \tarea{P1}{Probar las cuatro páginas, modal y carrusel en móvil, tablet y escritorio.}
  \tarea{P1}{Verificar navegación con teclado, orden y visibilidad de foco y activación de controles.}
  \tarea{P1}{Corregir elementos interactivos no semánticos cuando proceda.}
  \tarea{P1}{Probar movimiento reducido y asegurar que la animación no oculta contenido/controles.}
  \tarea{P2}{Clasificar imágenes decorativas/informativas y completar textos alternativos pertinentes.}
  \tarea{P2}{Guardar capturas comparables por ruta, dispositivo y estado.}
\end{listaacciones}
\criteriocierre{La web es legible y operable en los tamaños y modalidades aprobados, con evidencia de prueba.}

\chapter{Idiomas}

\begin{listaacciones}
  \tarea{P0}{Confirmar idiomas, alcance y responsable de traducciones.}
  \tarea{P1}{Elegir estrategia de internacionalización y trasladar textos a recursos localizados.}
  \tarea{P1}{Conectar el selector al idioma real de páginas, errores, proyectos y metadatos.}
  \tarea{P1}{Definir persistencia del idioma y su representación en la URL.}
  \tarea{P2}{Revisar formatos de fechas, etiquetas y textos traducidos.}
\end{listaacciones}
\criteriocierre{Los idiomas aprobados están completos o el selector se retira hasta poder ofrecerlos.}

\chapter{Analítica, errores y consentimiento}

\begin{listaacciones}
  \tarea{P0}{Decidir si se miden visitas y errores y con qué finalidad.}
  \tarea{P0}{Evaluar proveedor, datos, retención, ubicación y accesos con la cliente.}
  \tarea{P1}{Determinar con asesoramiento adecuado qué consentimiento y controles aplican.}
  \tarea{P1}{Si se aprueba, implementar medición con minimización de datos y sin datos sensibles.}
  \tarea{P1}{Probar aceptar, rechazar, retirar y persistir preferencias cuando aplique.}
  \tarea{P2}{Acordar responsable, umbrales y respuesta a alertas.}
\end{listaacciones}
\criteriocierre{Herramientas activas, avisos y controles coinciden; la decisión de no añadir analítica también queda registrada.}

\chapter{Recursos, caché y rendimiento}

\begin{listaacciones}
  \tarea{P1}{Medir cabeceras \texttt{Cache-Control}, \texttt{ETag} y \texttt{Last-Modified} en producción.}
  \tarea{P1}{Comparar primera y segunda visita; no deducir caché solo por ausencia de service worker.}
  \tarea{P1}{Decidir origen de imágenes/fuentes y comprobar permisos de uso.}
  \tarea{P1}{Bloquear Google Fonts y validar legibilidad, fallback y carga de fuente local.}
  \tarea{P2}{Medir peso/carga de imágenes y vídeo antes de optimizar.}
  \tarea{P2}{Acordar objetivos medibles de rendimiento.}
\end{listaacciones}
\criteriocierre{La estrategia se basa en mediciones del sitio publicado y se verifica en visitas repetidas y red limitada.}

\chapter{JavaScript, metadatos y PWA}

\begin{listaacciones}
  \tarea{P1}{Añadir explicación \texttt{\textless noscript\textgreater} y probar con JavaScript desactivado.}
  \tarea{P1}{Sustituir título y favicon de plantilla y corregir el atributo \texttt{lang}.}
  \tarea{P2}{Decidir si la instalación/PWA es requisito o queda fuera de alcance.}
  \tarea{P2}{Si se aprueba PWA, definir manifest, iconos, offline, service worker y actualizaciones.}
\end{listaacciones}
\criteriocierre{La experiencia sin scripts está explicada y el alcance PWA está aprobado o descartado explícitamente.}

\chapter{Calidad, seguridad y publicación}

\begin{listaacciones}
  \tarea{P1}{Ejecutar \texttt{npm run lint} y \texttt{npm run build}; registrar versión de Node.js, salida y advertencias.}
  \tarea{P1}{Acordar pruebas automatizadas; el proyecto no declara actualmente script de tests.}
  \tarea{P1}{Verificar navegadores y dispositivos acordados.}
  \tarea{P1}{Revisar proveedor, dominio, HTTPS, reglas SPA, cabeceras y proceso de despliegue.}
  \tarea{P1}{Realizar revisión de seguridad acorde con la arquitectura final y cualquier CMS/API/carga.}
  \tarea{P2}{Documentar copias, restauración, propiedad de cuentas y mantenimiento.}
  \tarea{P2}{Decidir si \archivo{analisis} se versiona: actualmente \archivo{.gitignore} excluye la carpeta; no cambiar la regla sin autorización.}
\end{listaacciones}
\criteriocierre{La entrega se valida en el entorno acordado y existe responsable de operación/recuperación.}

\end{document}
